\section{The intrinsic quotient}\label{quo:section}

We construct a quotient whose fibres are the projective-open varieties
found in the boundary windows. The first task is to show that these
fibres are intrinsic: two of them which meet must coincide. Guards
then make their translated charts exhaust the covering space, and
compactness of the inverse charts identifies the quotient with the
limiting tower. We begin by putting the given cover in the smooth
setting required by the window construction.

\subsection{The covering is locally pseudoconvex}

Normalize the projective ambient variety along the component containing
\(\widetilde X\), and use a functorial projective resolution
\cite[Theorem~1.1]{BierstoneMilman2008}.
Compatibility with local analytic isomorphisms identifies its restriction
with the pullback of a projective resolution \(X^+\to X\). Thus
\[
 \mu:V\longrightarrow\widetilde X
\]
is a proper modification, \(V\) is a smooth semialgebraic open subset
of a smooth projective variety \(Y\), and the lifted action of the
deck group \(\Gamma=\pi_1(X)\) is
free and properly discontinuous with compact projective quotient
\(X^+\). The inverse image \(V\) is connected: the fibres of a
proper modification of a normal space are connected, and such a map
cannot disconnect the inverse image of a connected open set.

\begin{lemma}\label{quo:pseudoconvex}
The open set \(V\) is locally pseudoconvex in \(Y\).
\end{lemma}
\begin{proof}
This is precisely Corollary~2(a) of Ivashkovich
\cite{Ivashkovich2008}. Here is also the Hartogs argument in the
present projective case. On a Hartogs figure in a coordinate polydisc
\(B\subset Y\), the covering map \(p:V\to X^+\) extends
meromorphically to \(B\), by the meromorphic Hartogs extension
theorem of Ivashkovich \cite{Ivashkovich1992}, in the form of
Theorem~4 of \cite{Ivashkovich2008}. It agrees with \(p\) on the
component of \(B\cap V\) containing the figure.

If that component does not fill \(B\), choose a straight segment
from a point of the figure to a point outside \(V\), and let \(a\)
be its first exit. The complex line containing the segment is not
contained in the indeterminacy locus, since it meets the figure.
Restriction of the meromorphic extension to this line extends
holomorphically across its isolated indeterminacies: in projective
homogeneous coordinates cancel the common vanishing order.
Consequently \(p\) tends to a point \(x\in X^+\) along the segment.
An evenly covered neighbourhood of \(x\) contains the image of a
connected final subsegment. That subsegment lies in one sheet and
converges, by the inverse sheet map, to a point of \(V\). Its limit
in the Hausdorff space \(Y\) must also be \(a\), a contradiction.
Every such Hartogs figure therefore fills in \(V\), which is the
local pseudoconvexity criterion. Pulling back its local
plurisubharmonic tests gives the same boundary property for
holomorphic inverse images used in the construction.
\end{proof}

\subsection{Whole projective fibres in a window}

The analytic-to-algebraic passages in this section use Remmert's
proper-image theorem in the reduced-space form proved in
\cite[\S7, Satz~1]{Grauert1960}, and Chow's theorem in
\cite[Proposition~13]{Serre1956}. We also use normalization and
normal holomorphic removability as in
\cite[II.7.3, II.7.12 and Remark~II.7.14]{Demailly2012}.
These apply to the possibly singular reduced incidence spaces as well
as to their smooth models.

We may now apply Theorem~\ref{win:construction}. Write \(n=\dim Y\),
\(r=n-d\), and denote the parameter family and its evaluation by
\[
 q:Q\longrightarrow B,\qquad e:Q\longrightarrow Y.
\]
The map \(q\) is smooth, proper, and projective, with connected
\(r\)-dimensional fibres. Their images \(Z_b=e(Q_b)\) are irreducible
projective varieties of dimension \(r\). The evaluation has generic
rank \(n\). The component \(S\) of \(q(e^{-1}(V))\) used in the
construction has a relative algebraic container for the complement
of \(G=e^{-1}(V)\) in \(Q|_S\), furnished by that theorem.
The same theorem makes the actual holes
\(J=(Q|_S)\setminus G\) analytic. The proper map
\((q,e):Q|_S\to S\times Y\) maps \(J\) to an analytic subset
of the incidence family, whose fibre at \(b\) is exactly
\(Z_b\setminus V\): every preimage of a point outside \(V\)
lies in \(J\), and no preimage of a point in \(V\) does.
Thus \(Z_b\cap V\) itself is a Zariski open subset of \(Z_b\),
not merely an image of some chosen Zariski open set.
All of these assertions concern whole fibres.

After normalization of the parameters, the accessible sets on successively
larger balls have marked open kernel \(P\). Every compact subset of
\(P\) eventually belongs to the accessible set. Their upper limits on
bounded sets belong to \(\overline P\). Boundary points other than those
on the artificial spheres lie outside \(q(e^{-1}(V))\). The last
assertion means that an incidence with a fixed point of \(V\) cannot
leave the chosen parameter component across a nonartificial boundary.
These are the interfaces with the construction and localization
sections; no conclusion about a quotient is part of these interfaces.

If \(d=0\), the window construction says that the whole of \(V\)
is a projective-open variety with analytic holes. Set \(M=P\) to
be a point and let \(f:V\to M\) be the constant map. The
normal-source assertions below follow from \(\mu\), and the
conditional proper-family conclusions concern maps to a point.
The guard, incidence, and inverse-map arguments that follow are
needed only when \(d>0\), which we assume until the descent step.

\subsection{The tower and its guards}

Write the successive symmetric matrices defining \(P\) as
\[
 M_j(Z)=
 \begin{pmatrix}s_j&w_j^{\mathsf T}\\w_j&\tau_j(Z_{>j})\end{pmatrix},
 \qquad j=1,\ldots,k.
\]
Here \(Z_{>j}\) consists of the later blocks, \(\tau_j\) is
holomorphic and polynomial, and every coordinate of \(Z\) occurs
as \(s_j\) or an entry of \(w_j\). The defining conditions are
\(\operatorname{Im}M_j(Z)>0\) for all \(j\).

\begin{lemma}\label{quo:tower}
The domain \(P\) is diffeomorphic to \(\R^{2d}\), is Stein, and is
biholomorphic to a bounded domain in \(\C^d\). Its Kobayashi distance
is a proper distance inducing its usual topology.
\end{lemma}
\begin{proof}
The first three assertions are Proposition~\ref{win:tower-properties}.
For the distance assertion, the map
\(Z\mapsto(M_1(Z),\ldots,M_k(Z))\) is a closed holomorphic
embedding into the product of the corresponding Siegel upper half
spaces. Closedness follows because the entries \(s_j,w_j\) recover
\(Z\), and a limit in that product still satisfies every strict
matrix inequality. The Cayley transform identifies each Siegel
upper half space with the unit ball for the operator norm in
symmetric matrices. The Kobayashi distance from its origin is
\(\operatorname{arctanh}\|W\|\): the radial disc gives the upper
bound, and a norming complex linear functional gives the lower
bound. Its automorphisms act transitively, so its distance balls
are compact. Distance decrease therefore puts a Kobayashi
ball of \(P\) inside a compact subset of \(P\). Bounded realization
implies nondegeneracy, and coordinate balls give the local upper
estimates proving that the Kobayashi distance induces the usual
topology. The balls are consequently closed in those compact sets,
and hence compact.
\end{proof}

\begin{lemma}\label{quo:guard}
There are normalized accessible parameter regions \(A_i\subset\{|Z|<R_i\}\),
where \(R_i\to\infty\), and holomorphic functions \(h\) on their
neighbourhoods, given by the same formula
\begin{equation}\label{quo:guard-formula}
 h(Z)=\exp\left(-\alpha\sum_j\operatorname{tr}
                      \log(I-iM_j(Z))\right),
\end{equation}
with the following properties. Their values lie in a fixed sector
\(\{\zeta:|\arg\zeta|<\theta\}\), \(\theta<\pi/2\), and are
uniformly bounded. On \(A_i\), \(|h(Z)|\leq C(1+|Z|)^{-\alpha}\).
Every compact subset of \(P\) is eventually contained in \(A_i\).
There are connected open subsets \(H_i\subset A_i\), containing
the seed and eventually every compact subset of \(P\), such that
\[
 u_i=\log|h|+a_i|Z|^2>m_i\quad\hbox{on }H_i,
 \qquad a_i>0,
\]
where \(a_iR_i^2\leq1\), \(m_i\to-\infty\), and every boundary
of \(H_i\) admitting an incidence with a point of \(V\) has
\(|h|\leq\delta_i\), with \(\delta_i\to0\).
\end{lemma}
\begin{proof}
Choose the scale on the ball of radius \(R_i+1\) so that localization
gives \(\operatorname{Im}M_j\geq-\eta_i I\), with
\(\eta_i\downarrow0\) and \(\eta_i<1/2\), throughout the accessible
component. Require also the distance of its points to
\(\overline P\), on each previously chosen bounded ball, to tend
to zero. This is a diagonal use of fixed-ball estimates.

For a matrix with \(\operatorname{Re}A_j\geq I/2\), the principal
matrix logarithm is holomorphic and
\(|\operatorname{Im}\operatorname{tr}\log A_j|<n_j\pi/2\),
where \(n_j\) is its size. Choose
\(0<\alpha<1/(2\sum_jn_j)\). This proves the sector bound.
All singular values of \(A_j\) are at least \(1/2\); hence
\[
 |\det A_j|\geq2^{-n_j},\qquad
 |\det A_j|\geq2^{-(n_j-1)}\|A_j\|.
\]
Since the matrices contain every coordinate of \(Z\), multiplying
these inequalities proves the asserted decay estimate.

Let \(r_i\downarrow0\) bound \(|h|\) near the artificial sphere.
Take \(a_iR_i^2\leq1\), and choose \(m_i\to-\infty\) with
\(\log r_i+1<m_i\). Let \(H_i\) be the seed component of
\(\{u_i>m_i\}\) in \(A_i\). Any connected compact subset of
\(P\) containing the seed belongs to \(H_i\) eventually, since
\(h\) has a positive minimum modulus there. At a level boundary
\(|h|\leq e^{m_i}\); the artificial boundary lies below the level;
and the other boundaries admit no incidence with \(V\). Take
\(\delta_i=e^{m_i}\). The function \(u_i\) is strictly
plurisubharmonic because \(h\) is nowhere zero.
\end{proof}

\subsection{Incidence traces and whole fibres}

After shrinking \(B\) at a generic real centre, the cycle map
\(b\mapsto Z_b\) is injective. To justify this, if its generic
rank were less than \(d\), the union of its \(r\)-dimensional
cycles would have dimension less than \(r+d=n\), contradicting
the generic rank of evaluation. The constant-rank theorem then
gives an injective restriction. Let
\[
 \mathcal I=\{(b,y)\in B\times Y:y\in Z_b\}.
\]
We give it its reduced structure; it is a proper analytic family
over \(B\), of pure dimension \(n\), obtained as the proper image
of \(Q\). The following argument is made separately at each fixed
scale; its bounds need not be uniform in the scale.

\begin{lemma}\label{quo:incidence}
The union \(U_i=\bigcup_{b\in H_i}(Z_b\cap V)\) is open in \(V\).
Distinct fibres in this union are disjoint, and there is a
holomorphic submersion
\[
 \pi_i:U_i\longrightarrow H_i,
 \qquad \pi_i^{-1}(b)=Z_b\cap V.
\]
These fibres are connected. Moreover, if \(T\) is a smooth connected
projective variety, \(D\subset T\) is a proper analytic subset,
and \(g:T\setminus D\to V\) is holomorphic, then any such image
meeting a fibre of \(\pi_i\) is contained in that fibre.
\end{lemma}
\begin{proof}
Fix \(y\in V\). The set of its incident parameters is analytic.
A positive-dimensional irreducible component meeting \(H_i\)
would have a maximum of \(u_i\) greater than \(m_i\): on the
artificial boundary \(u_i<m_i\), while the other parameter
boundaries cannot contain an incidence with \(y\). The part above
any intermediate level between \(m_i\) and a value on that component
is compact. The maximum principle for strictly plurisubharmonic
functions on positive-dimensional analytic sets is a contradiction.
Thus all incident parameters in \(H_i\) are isolated.

Near any one of these roots, the projection of \(\mathcal I\) to
\(V\) is proper and finite after choosing a small parameter
neighbourhood whose boundary avoids the root. It is open: each
local irreducible component of the finite source has dimension
\(n\), so its image germ is the full germ of the smooth
\(n\)-dimensional target. Count the roots with their finite-map
multiplicities. The function
\begin{equation}\label{quo:trace}
 \Phi_i(y)=\sum_{b:y\in Z_b,\ b\in H_i}
                    \max(0,u_i(b)-m_i)
\end{equation}
is bounded and plurisubharmonic on \(V\), with value zero when
there are no positive summands. Here and below one may first use
a slightly larger high region and put the cutoff strictly inside
it. This avoids any finiteness assertion exactly on a level seam.
Away from a finite projection's discriminant the assertion is the
sum rule for plurisubharmonic functions on holomorphic branches.
Across the discriminant the roots remain in a compact parameter
set, so the bounded removable extension gives the same trace,
including multiplicities. Across a cutoff seam the summands tend
to zero, and the plurisubharmonic pasting lemma applies. No upper
part can appear without a limiting incidence. Finally, the
subanalytic incidence family has a uniform finite bound on its
number of isolated roots on this fixed ball. Multiplicity is
bounded by the corresponding number of simple roots at nearby
generic target points. This proves boundedness of the trace.

A bounded-above plurisubharmonic function on a Zariski open subset
of a connected smooth projective variety extends across its
analytic complement by \cite[I.5.24]{Demailly2012} and is constant
by compactness. In particular
\(\Phi_i\) is constant after pullback to the projective model
of any \(Z_b\cap V\), \(b\in S\), or by any map \(g\) in the
statement. On a small disc in that model, lift the finite
correspondence of high roots to its normalization and make a
finite ramified base change. All its branches then become
holomorphic. The trace is a sum of subharmonic functions, and
each high branch contributes \(u_i\) minus a constant. Since
the sum is constant and \(u_i\) is strictly plurisubharmonic,
each high parameter branch has zero derivative. This also proves
constancy at a multiple root, by continuity.

If high fibres \(Z_b\cap V\) and \(Z_c\cap V\) meet, apply this
observation on a projective model of the first fibre. A neighbourhood
of the meeting point in that fibre lies in \(Z_c\). Irreducibility
and equal dimension give \(Z_b=Z_c\), hence \(b=c\). The finite
incidence projection consequently has one point in each fibre
over \(U_i\). Its reduced graph is a finite bimeromorphic map
onto a normal space, so it is an isomorphism. This proves
holomorphy and openness. Since \(\pi_i\circ e=q\) over the high
region and \(q\) is a submersion, \(\pi_i\) is a submersion at
every point: choose any preimage in \(Q_b\). The complement of
a proper analytic subset of the connected smooth \(Q_b\) is
connected and maps onto \(Z_b\cap V\); hence that fibre is
connected. Finally the argument with \(g\) gives a constant
parameter near a meeting point. The inverse image of the closed
analytic subset \(Z_b\cap V\) then contains an open set in the
connected source \(T\setminus D\), so is the whole source.
\end{proof}

The induced guard \(h\circ\pi_i\) has the boundary estimate
\begin{equation}\label{quo:chart-boundary}
 \limsup_{U_i\ni y'\to y}|h(\pi_i(y'))|\leq\delta_i
 \qquad(y\in\partial U_i\cap V).
\end{equation}
Indeed any offending sequence has a subsequence whose parameters
converge in the closed artificial ball. Closedness of incidence
gives an incidence with \(y\). The limit cannot belong to \(H_i\),
which would put \(y\) in \(U_i\), and cannot lie on a
nonartificial boundary outside \(q(G)\). It is therefore one of
the low-guard boundaries in Lemma~\ref{quo:guard}. This proves
the estimate without a bound on tangent directions along the
fibres.

\subsection{A compactness lemma for cutoff guards}

The next lemma allows the domains of the guards themselves to move.
The small-boundary assumption is uniform: it applies at every
boundary point internal to the test manifold.

\begin{lemma}\label{quo:cutoff-normal}
Let \(D\) be a complex manifold, let \(A_i\subset D\) be open,
and let \(h_i\in\mathcal O(A_i)\) satisfy \(|h_i|\leq C\).
Suppose that, for every \(x\in\partial A_i\cap D\),
\[
 \limsup_{A_i\ni z\to x}|h_i(z)|\leq\epsilon_i,
 \qquad \epsilon_i\longrightarrow0.
\]
Extend \(h_i\) by zero outside \(A_i\). These extended functions
have a subsequence converging uniformly on compact subsets of
\(D\) to a holomorphic function. Suppose in addition that their
nonzero values lie in \(|\arg\zeta|\leq\theta<\pi/2\), that
\(D\) is connected, and that
\(x_i\to x\), \(h_i(x_i)\to a\ne0\). Then every resulting
limit is nowhere zero, and every compact subset of \(D\) lies
in \(A_i\) eventually.
\end{lemma}
\begin{proof}
Increase \(\epsilon_i\) by \(1/i\) if necessary so it is positive.
First work on a disc compactly contained in \(D\), enlarging it
slightly for estimates. Put \(\eta_i=2\epsilon_i\), and extend
\(v_i=\max(|h_i|,\eta_i)\) by the constant \(\eta_i\).
Pasting gives a uniformly bounded subharmonic function. If
\(\chi\) is a smooth cutoff equal to one on a smaller disc,
\(\int\chi\,\Delta v_i=\int v_i\Delta\chi\) gives a uniform
bound for its Riesz mass there. On \(|h_i|>\eta_i\),
\[
 \Delta|h_i|=\frac{|h_i'|^2}{|h_i|}
\]
with the usual Euclidean Laplacian. It follows that the energy
of \(h_i\) on this set is bounded, and that
\[
 \int_{2\eta_i<|h_i|<4\eta_i}|h_i'|^2=O(\eta_i).
\]
Choose a smooth scalar cutoff \(\lambda\), zero on \([0,2]\)
and one on \([4,\infty)\), and extend
\(F_i=\lambda(|h_i|/\eta_i)h_i\) by zero. The extension is
locally in \(W^{1,2}\): near every internal boundary the formula
vanishes, and the preceding energy bounds control the remaining
region. Moreover
\[
 \|F_i\|_{W^{1,2}(D')}\leq C_{D'},\qquad
 \|\bar\partial F_i\|_{L^2(D')}^2=O(\eta_i),\qquad
 |F_i-h_i|\leq4\eta_i
\]
for the zero extension in the last expression. Rellich compactness
and the distributional Cauchy--Riemann equation show that a
subsequence converges in \(L^2_{\rm loc}\) to a holomorphic
function \(h\).

We verify uniform convergence, which does not follow just from
this Sobolev bound. Fix an annulus compactly inside the test disc
and first pass to a subsequence for which the squared \(L^2\)
errors \(\|F_i-h\|^2\) on that annulus are summable. Fubini's
theorem gives \(L^2\) convergence of the traces for almost every
radius. The integrals in the radial variable of the tangential
\(W^{1,2}\) energies are uniformly bounded. Fatou's lemma shows
that their lower limit is finite for almost every radius. Choose
a radius satisfying both conclusions and then a further subsequence
with bounded tangential energies on that circle. Compactness of
\(W^{1,2}\) of a circle in the continuous functions gives uniform
trace convergence; the difference between \(F_i\) and the zero
extension of \(h_i\) tends uniformly to zero as well.
If \(h(z_0)\ne0\), start with an arbitrarily small annulus about
\(z_0\) on which \(h\) takes values in a disc \(B(a,r)\) with
\(0\notin\overline{B(a,2r)}\). The preceding radius choice then
supplies such a surrounding circle inside this annulus.
For large \(i\), the original \(h_i\) is defined near that
circle and its values belong to \(B(a,2r)\).

For completeness, the argument principle excludes hidden holes
inside this circle. Choose a regular level
\(t_i\in(4\eta_i,5\eta_i)\) and apply the argument principle
to the part of \(\{|h_i|>t_i\}\) inside the circle. This part
has finitely many boundary components after a harmless choice
of the level and outer circle: it is compact in \(A_i\), and
the level is real analytic and regular. On each inner boundary,
the image has modulus \(t_i\), so its winding number about a
value \(w\) with \(|w|>t_i\) is zero. For
\(w\notin\overline{B(a,2r)}\), the outer winding number is
also zero. Thus \(h_i\) cannot take such a value with
\(|w|>t_i\) inside. The component adjoining the outer circle
cannot reach an inner level boundary: a continuous image from
\(B(a,2r)\) to that small circle would pass through an excluded
value. It therefore fills the disc, and all its values belong
to \(\overline{B(a,2r)}\). Taking arbitrarily small circles
and radii proves locally uniform convergence near \(z_0\).

At a zero of \(h\), choose a surrounding circle with no zeros
of \(h\) and small maximum modulus, and apply the maximum
principle to the pasted \(v_i\). This proves uniform convergence
also there. If \(h\equiv0\), the sliced circles and the same
maximum principle suffice directly. Hence every sequence has
a uniformly convergent subsequence on a smaller disc.

The estimates were uniform for discs of fixed radius lying in
a fixed coordinate box. They consequently apply to moving
complex lines as well. They imply \emph{asymptotic}
equicontinuity of the zero extensions: if it failed, one could
choose indices tending to infinity and two converging points
with the same limit and a fixed positive difference of values. Apply the
one-variable compactness result on the complex lines through
these pairs, using fixed-radius parametrized discs in the box,
to obtain a contradiction. The zero extensions can still have
small jumps at each fixed index, so apply Arzela--Ascoli to the
continuous functions \(F_i=\lambda(|h_i|/\eta_i)h_i\), defined
by the same cutoff formula and by zero in all dimensions.
Their difference from the zero extensions is at most
\(4\eta_i\). Asymptotic equicontinuity, together with continuity
of each of the finitely many initial \(F_i\), gives ordinary
equicontinuity of this family on compact subsets. Arzela--Ascoli
and a countable exhaustion give local uniform convergence of
\(F_i\), and hence of the zero extensions. The limit is
holomorphic on every coordinate line, by the disc result,
hence holomorphic.

Finally \(\operatorname{Re}h\geq0\), and the moving-point
condition gives \(h(x)=a\ne0\). The strong minimum principle
gives \(\operatorname{Re}h>0\) on connected \(D\). Its modulus
has a positive minimum on every compact set. Uniform convergence
then excludes points where the zero extension was used.
If compact containment failed for the original sequence, a
violating subsequence would have a further uniformly convergent
subsequence with the same nowhere-zero conclusion. This proves
eventual compact containment for the original sequence as well.
\end{proof}

\begin{corollary}\label{quo:trapping}
Let \(D_i\) exhaust a connected complex manifold \(D\), and let
\(g_i:D_i\to V\) be holomorphic. Suppose that at points
\(x_i\to x\) the images belong to \(U_i\), and their normalized
parameters \(\pi_i(g_i(x_i))\) converge to the seed \(p\in P\).
Then every compact subset of \(D\) is eventually mapped into
\(U_i\); after taking a subsequence, \(\pi_i\circ g_i\) converges
normally on \(D\) to a holomorphic map into \(P\). The same
assertion holds for maps into a manifold equipped with these
projection charts and their guards.
\end{corollary}
\begin{proof}
On each connected relatively compact test domain containing
\(x\), the maps \(g_i\) are defined for all sufficiently large
\(i\). Apply Lemma~\ref{quo:cutoff-normal} there to
\(h\circ\pi_i\circ g_i\) on \(g_i^{-1}(U_i)\).
The boundary estimate is \eqref{quo:chart-boundary}, including
where the projection chart ceases to exist. At the moving seed
the guard tends to \(h(p)\ne0\). A diagonal exhaustion of \(D\)
therefore gives eventual inclusion in \(U_i\) and a nowhere-zero
guard limit. Its positive minimum modulus on compact sets,
together with the decay in Lemma~\ref{quo:guard}, bounds all
normalized parameters there. Montel's theorem gives a holomorphic
parameter limit \(F\).

To apply localization, join \(x\) to a prescribed point of \(D\)
by a compact path, and precede it by a short path from \(x_i\)
to \(x\) in a fixed coordinate ball. Include these paths in one
relatively compact test domain. The preceding bounds hold
on the whole paths, whose images eventually lie in \(U_i\).
Their parameter paths start at points tending to \(p\) and stay
in one fixed bounded box. The moving-start assertion of
Theorem~\ref{loc:theorem} puts their limiting endpoints in
\(\overline P\). Hence \(F(D)\subset\overline P\). For each constant vector \(v\),
\(v^*\operatorname{Im}M_j(F)v\) is a nonnegative harmonic
function, positive at the seed whenever \(v\ne0\). The minimum
principle makes it positive everywhere. Thus all the strict
matrix inequalities hold and \(F(D)\subset P\).
\end{proof}

\subsection{Exhaustion and the inverse maps}

Choose a point \(y_i\) over the normalized seed of \(\pi_i\).
Cocompactness supplies \(\gamma_i\in\Gamma\) and, after passing
to a subsequence, points \(q_i\to q\in V\) with
\(\gamma_iq_i=y_i\). Corollary~\ref{quo:trapping}, applied to
\(\gamma_i:V\to V\), proves that
\[
 \widehat U_i=\gamma_i^{-1}U_i
\]
eventually contains every compact subset of \(V\). In particular
these sets form a covering, although they need not be nested.
Write \(\widehat\pi_i=\pi_i\circ\gamma_i\).

Any two fibres of these charts, or of any of their deck translates,
which meet are equal. Indeed each has a smooth projective model
with an analytic complement, and the maximality assertion of
Lemma~\ref{quo:incidence} gives both inclusions. They therefore
partition \(V\) into connected entire fibres. Give their set
\(M\) the quotient topology. Each \(\widehat U_i\) is saturated
and its quotient is biholomorphic to \(H_i\). On overlaps the
coordinate change is holomorphic: use a local holomorphic
section of the submersion \(\widehat\pi_i\) and compose with
the other projection. Equality of whole fibres makes the result
independent of the section. The inverse change is constructed
in the same way.

The quotient is Hausdorff. Given two distinct leaves, choose one
point on each. They lie together in one sufficiently late
\(\widehat U_i\), which contains both entire leaves. Disjoint
parameter neighbourhoods in that chart have disjoint saturated
inverse images, open also in \(V\). The quotient is second
countable because the quotient map is open and \(V\) is second
countable. It is connected and is therefore a complex manifold.
We obtain a surjective submersion
\[
 f:V\longrightarrow M.
\]
The deck action descends to \(M\).
Let \(M_i=f(\widehat U_i)\), let
\(\varphi_i:M_i\xrightarrow{\sim}H_i\) be its coordinate map,
and let \(\psi_i:H_i\to M_i\) be its inverse.
In particular, for any connected parameter box \(B_0\Subset H_i\),
the translated evaluation from \(G\cap q^{-1}(B_0)\) maps onto
the \emph{entire} saturated set \(f^{-1}(\varphi_i^{-1}(B_0))\).
It is generically of full rank and comes from the smooth proper
projective family \(Q|_{B_0}\). Its actual omitted set is analytic
and is contained in the specified relative algebraic divisor.
The next section applies meromorphic extension to this proper
family and uses surjectivity onto each entire leaf.

Every compact subset of \(M\) eventually belongs to \(M_i\): finitely many
local sections of \(f\) lift that compact set to a compact
subset of \(V\). Corollary~\ref{quo:trapping} and local sections
give a normally convergent subsequence
\[
 \varphi_i\longrightarrow\varphi:M\longrightarrow P.
\]
If \(b_i=f(q_i)\) and \(b_0=f(q)\), then
\(b_i\to b_0\), \(\varphi_i(b_i)=p\), and
\(\varphi_i(b_0)\to p\).

\begin{lemma}\label{quo:inverse-tightness}
For every compact \(K\subset P\), the images \(\psi_i(K)\)
eventually lie in a compact subset of one fixed chart \(M_j\).
\end{lemma}
\begin{proof}
Enlarge \(K\) to a compact set contained in a connected relatively
compact domain of \(P\) containing \(p\). Suppose the assertion
fails. Choose \(j_\nu\to\infty\) and compact sets
\(L_\nu\subset H_{j_\nu}\cap P\) exhausting \(P\): each
fixed compact subset of \(P\) belongs to \(L_\nu\) eventually.
Failure means that, for each fixed \(j_\nu,L_\nu\), there are
arbitrarily large \(i\) for which
\(\psi_i(K)\not\subset\varphi_{j_\nu}^{-1}(L_\nu)\).
Choose one such \(i_\nu\) so large that \(H_{i_\nu}\) contains
the \(\nu\)-th member of a fixed exhaustion of \(P\), that
\(b_{i_\nu}\in M_{j_\nu}\), and that
\[
 |\varphi_{j_\nu}(b_{i_\nu})-
   \varphi_{j_\nu}(b_0)|<1/\nu.
\]
These last conditions are possible because \(j_\nu\) is fixed
when \(i_\nu\) is chosen. The displayed starting parameters
therefore tend to \(p\).

Test \(\psi_{i_\nu}\) against the chart \(M_{j_\nu}\) and its
guard. Corollary~\ref{quo:trapping} applies on the expanding
domains in \(P\). It says that \(\psi_{i_\nu}(K)\subset
M_{j_\nu}\) eventually, and that their tested parameters lie
in one fixed compact subset \(L\subset P\), after a subsequence.
For large \(\nu\), \(L\subset L_\nu\), contradicting the
choice of \(i_\nu\). This proves the fixed-chart assertion;
control in merely changing charts would not have sufficed.
\end{proof}

By this lemma and Montel in the fixed charts, a diagonal
subsequence of \(\psi_i\) converges normally to
\(\psi:P\to M\). Passing to the limits in
\(\varphi_i\psi_i=\id\) and \(\psi_i\varphi_i=\id\)
is legitimate on compact sets, by both normal convergences.
Thus \(\varphi\) and \(\psi\) are inverse biholomorphisms.
In particular \(M\simeq P\), so the complete real holomorphic
vector fields of Proposition~\ref{win:tower-properties} transfer
to \(M\) and span every tangent space over \(\C\).

\subsection{Normal-source descent and regularity}

\begin{lemma}\label{quo:descent}
The map \(f\) descends uniquely to an equivariant holomorphic
map \(f_X:\widetilde X\to M\). Its fibres are connected and,
locally on \(M\), have simultaneous projective compactifications
with analytic holes. Each fibre is biholomorphic to a semialgebraic
open subset of a projective variety.
\end{lemma}
\begin{proof}
Realize \(M\) as a bounded domain. Each coordinate of \(f\)
is constant on every compact connected fibre of \(\mu\): a
holomorphic function on a compact irreducible complex space is
constant, and connectedness makes the constants on its components
agree. For completeness, a holomorphic function on the inverse image
of a small open set under a proper modification descends on its
isomorphism locus and is locally bounded by properness. Normal
removability extends it across the omitted analytic subset.
Consequently \(\mu_*\mathcal O_V=\mathcal O_{\widetilde X}\),
which gives the unique holomorphic descent. Its values lie in \(M\), and uniqueness
gives equivariance. Also
\(\mu^{-1}(f_X^{-1}(m))=f^{-1}(m)\), which proves connectedness.

On an original high chart, apply the projective ambient resolution
map to the compact incidence family. Its proper image in parameter
times the normal ambient compactification is analytic and proper
over parameters. Its restriction to \(\widetilde X\) is exactly
the graph of the single-valued descended parameter: proper
surjectivity of \(\mu\) and the preceding fibre identity prove
both inclusions. Normalize this compactifying family if necessary;
normalization does not change the normal graph portion. The graph
portion is open, and its fibre complement is analytic and hence
algebraic in each projective fibre. Equivalently, before a deck
translation each fibre is the intersection of an algebraic image
cycle with the original semialgebraic \(\widetilde X\). This also
proves the semialgebraic assertion. Translated charts preserve
the biholomorphic assertion, without requiring algebraic deck
transformations.
\end{proof}

\subsection{Preparing the normal compactifying family}

To use the normal-source fibres in the Albanese construction, we
need normality and flatness at every parameter, a resolution on
every fibre, and topological local triviality once the deck action
on the base is free and discrete. All three properties will come
from the original compact incidence families. We have so far
described the quotient for a window at generic real centres.
Before making the final choice of centre, scales, and deck
translations, impose the following conditions on its full
parameter box. The window construction permits these generic
shrinkings, and the preceding quotient arguments then apply to
the resulting choices.

Let \(\mathcal N\to B\) be the normalization of the proper
image of \(\mathcal I\subset B\times Y\) in parameter times
the normal ambient compactification, as in the descent proof.
Its open graph portion is the corresponding open subset of
\(\widetilde X\), and normalization leaves that portion unchanged.
The source to be resolved is the reduced incidence space
\(\mathcal I\), rather than its parametrizing space \(Q\).
Lemma~\ref{quo:incidence} identifies its entire high graph
portion with an open subset of \(V\), so this portion is smooth.
On these graph portions the ambient resolution induces exactly
\(\mu\).

\paragraph{Normal fibres and flatness.}
Shrink the full box so that \(\mathcal N\to B\) is flat with
normal fibres. The analytic result giving this shrinking is
\cite[Theorem~6.2(2)(ii), including the properness clause]{Greuel2017}:
for a proper morphism from a normal complex space to a smooth
complex space, the image of the locus where the morphism is not
normal is a nowhere dense closed analytic subset. Here a normal
morphism means flat with normal fibres. Thus this theorem applies
directly to \(\mathcal N\) over the polydisc;
no algebraization of the analytic coefficient functions is needed.
A proper complex analytic subset cannot contain an open part of
the full-dimensional totally real box, so a generic real centre
and a neighbourhood of it avoid that image.

\paragraph{A resolution on every fibre.}
The correspondence from \(\mathcal I\) to \(\mathcal N\)
induced by the ambient resolution is bimeromorphic: the full-rank
evaluation meets the locus where that resolution is an isomorphism.
Resolve its source over the full box, leaving the smooth graph
portion unchanged. Normality of the resolved source lifts the
map to the normalized target, giving a proper projective map
\[
                 \mathcal R\longrightarrow\mathcal N
                         \quad\hbox{over }B.
\]
Shrink generically so that \(\mathcal R\to B\) is smooth.
The exceptional sets of the map and its inverse birational
correspondence have smaller dimension than the total spaces.
The fibre-dimension theorem and properness allow another generic
shrinking so their intersections with every fibre have smaller
dimension than that fibre. The map on every fibre is consequently
bimeromorphic. On the graph portions it is precisely the original
map \(\mu\).

\paragraph{Submersivity on the intrinsic strata.}
Give the absolute space \(\mathcal N\) its canonical minimal complex Whitney
stratification. This stratification is intrinsic and restricts
to open subsets; this is Teissier's canonical complex-analytic Whitney
stratification \cite[VI, \S3, Propositions~3.1--3.2]{Teissier1982};
see also \cite[\S4.3, following Theorem~4.13]{GilesFloresTeissier2018}. After a further generic
shrinking, the parameter projection is a submersion on every
stratum. Here is why the shrinking can be made simultaneously.
Over a relatively compact parameter box, properness leaves only
finitely many strata to consider. The closure of their critical
loci is stratified-critical: if critical tangent spaces approach
a lower stratum, Whitney condition (a) puts its tangent space
inside their limit, and the limiting rank is still less than
the base dimension. Their image is closed by properness. It is
a subanalytic set of real dimension less than \(2d\), by the
rank theorem on the complex analytic strata. In fact the critical
image is complex analytic, by the proper-image theorem applied
to the analytic closures of the critical loci. To make this last
point precise, the canonical complex stratification has analytic
stratum closures and frontiers. Form the Nash modification of
the analytic closure of each stratum, which records the limiting
tangent planes in a Grassmannian bundle. The rank-drop condition
on its tautological tangent bundle is analytic, and the Nash
projection is proper. Whitney condition (a) gives the asserted
containment at lower strata after projecting back.
Thus the critical image's intersection
with the totally real box has empty interior. Remove that image.

\paragraph{From the prepared family to every fibre.}
Make these finitely many shrinkings before choosing the diagonal
scales in Lemma~\ref{quo:guard} and the deck translations used
for exhaustion. On each high chart, the graph portions of
\(\mathcal R\to\mathcal N\) now give the flat normal family,
its fibrewise resolution, and the stratum-submersion property.
Flatness and normality are invariant under local analytic
isomorphisms, as is the absolute Whitney stratification. Deck
translation therefore preserves all these properties.
The translated high charts exhaust the source and contain
entire fibres. Consequently the properties hold at
\emph{every} fibre and point of \(f_X\), not only over a
generic part of its base. The map \(f\) on \(V\) was already
a submersion; its induced maps after proper quotient are smooth
families. The same exhaustion makes \(\mu\) bimeromorphic on
every fibre.

The compactifications used to verify these assertions do not
replace the resolution on overlaps.
Write \(U=\widetilde X\) in the following equivariance identity.
The lifted deck transformations satisfy
\(\mu\circ\gamma_V=\gamma_U\circ\mu\), so each translated graph
portion still carries the same global map \(\mu:V\to U\).
In particular, for every subgroup \(K'\) used in a finite covering
family below, the resulting resolution morphism is the actual
quotient map \(V/K'\to U/K'\). Its maps on first cohomology are
induced by this single proper morphism, independently of the
incidence chart used to verify fibrewise regularity.

\begin{theorem}[Intrinsic projection]\label{quo:projection}
Under the hypotheses of condition~\textup{(i)} of
Theorem~\ref{thm:main} and the boundary construction of
Theorem~\ref{win:construction}, there are a contractible Stein
manifold \(M\), biholomorphic to the boundedly realizable tower
\(P\), an equivariant surjective submersion \(f:V\to M\), and
an equivariant holomorphic map \(f_X:\widetilde X\to M\) with
\(f=f_X\mu\). Both maps have connected fibres. Their fibres
have simultaneous projective compactifications with analytic
holes locally on \(M\), after a deck translation. The
normal-source fibres are normal and semialgebraically
compactifiable up to biholomorphism, and \(f_X\) is flat.

Suppose additionally that the image \(\Lambda\) of the deck
action on \(M\) is discrete and acts freely. Then the descended
proper map \(X\to M/\Lambda\), its pullback to \(M\), and
every connected finite covering family of that pullback are
flat and topologically locally trivial, with connected normal
projective fibres. Their induced resolution families are smooth
proper projective families; the resolution maps are proper,
have connected point fibres, and are bimeromorphic on every
fibre. The induced maps on integral degree-one cohomology are
morphisms of local systems.
\end{theorem}
\begin{proof}
Only the last paragraph remains. A discrete group acting on a
bounded domain acts properly; freeness makes the quotient map
a covering. Equivariance descends \(f_X\) to a holomorphic map
from compact \(X\), hence a proper map. Locally it has the
normal graph descriptions above. Flatness and normal fibres
follow from those descriptions. The intrinsic Whitney strata
downstairs are submersive over the base, since this is true on
the covering charts. Thom's first isotopy lemma, in the exact
proper stratified-submersion form of
\cite[Proposition~11.1]{Mather1970}, gives topological local
triviality. Its fibres are connected because upstairs fibres
are connected, and projective because they are compact analytic
subvarieties of projective \(X\).

Base change to \(M\) preserves these conclusions. A finite
covering of its total space is locally analytically isomorphic
to it; hence flatness, normality, and stratified submersivity
persist, and properness follows from finiteness. Apply the same
isotopy lemma. Its fibres are connected if the covering total
space is connected: the component covering of the simply
connected base \(M\) is trivial. The fibres are finite covers
of projective varieties and are projective. The resolution
families are proper submersions, so Ehresmann's theorem makes
their integral cohomology locally constant. The normal families
are topologically locally trivial by the preceding argument;
functoriality of cohomology for the proper fibre maps gives the
asserted morphism of local systems. Connectedness of the point
fibres and fibrewise bimeromorphicity follow from \(\mu\) and
the generic preparation already propagated through all charts.
\end{proof}
